%% ma: L12-B08-0131
%% lop: 12
%% chuong: 2
%% bai: 8
%% dang: 12.8.22
%% loai: TN4
%% muc: TH
%% dap_an: "B"
%% nguon: "CĐ2.3.pdf giáo viên gửi 27/10/2026 (M07, Lê Minh Khôi – Nguyễn Anh Khôi, Chương II Chủ đề 3 Biểu thức tọa độ của các phép toán vectơ), Câu 46"
%% kien_thuc: "Bài 8 (trung điểm, tích vô hướng)"
%% trang_thai: chinh-thuc
\begin{ex}
Trong không gian $Oxyz$, cho hai điểm $A(1;2;3)$ và $B(-1;4;1)$. Gọi $I$ là trung điểm của đoạn thẳng $AB$. Tính giá trị $\overrightarrow{IA}\cdot\overrightarrow{BA}$.
\choice{$\overrightarrow{IA}\cdot\overrightarrow{BA}=4$}{\True $\overrightarrow{IA}\cdot\overrightarrow{BA}=6$}{$\overrightarrow{IA}\cdot\overrightarrow{BA}=-4$}{$\overrightarrow{IA}\cdot\overrightarrow{BA}=-6$}
\loigiai{$I(0;3;2)$, $\overrightarrow{IA}=(1;-1;1)$, $\overrightarrow{BA}=(2;-2;2)$, tích vô hướng bằng $2+2+2=6$.}
\end{ex}
